\documentclass{article}
\usepackage[UTF8]{ctex}
\usepackage{amsmath,amssymb,amsfonts}
\usepackage{enumitem}
\usepackage{geometry}
\geometry{a4paper, margin=1in}
% === 错题本排版优化（不参与题目标识） ===
\setlength{\parindent}{0pt}
\setlength{\parskip}{0.35em plus 0.10em minus 0.08em}
\setlength{\emergencystretch}{3em}
\setlist{itemsep=0.25em, topsep=0.40em, parsep=0pt}
\allowdisplaybreaks[2]
% === 错题本排版优化结束 ===
\begin{document}
\section*{苏州大学2024年数学分析真题}
\section*{88．苏州大学2024年数学分析试卷}

\subsection{题目1}
1．求极限 \(\lim _{x \rightarrow 0} \frac{\sin x-\arctan x}{x^{3}}\) ．

2．设 \(f(x)=x^{3} \cos x\) ，求 \(f^{(2023)}(0)\) ．

\subsection{题目2}
二．（15 分）设 \(a_{1}=1\) ，且当 \(n \geq 2\) 时，\(a_{n}=\frac{1}{2}\left(a_{n-1}+\frac{5}{a_{n-1}}\right)\) 。证明：\(\left\{a_{n}\right\}\) 收敛，并求极限。

\subsection{题目3}
三．（15 分）设 \(f(x)=\left\{\begin{array}{ll}\frac{\sin x}{x}, & x \neq 0 ; \\ 1, & x=0 .\end{array}\right.\) 证明：函数 \(f(x)\) 在 \((-\infty,+\infty)\) 上一致连续．

\subsection{题目4}
四．（15 分）计算三重积分 \(\iiint_{\Omega}(x+y+z)^{2} \mathrm{~d} x \mathrm{~d} y \mathrm{~d} z\) ，其中 \(\Omega=\left\{(x,y,z)\in\mathbb{R}^{3}\mid |x|+|y|+|z|\leq1\right\}\) ．

\subsection{题目5}
五．（15 分）设 \(u=f(x, y, z)\) 的所有二阶偏导数都连续，\(v=f(x(s, t, r), y(s, t, r), z(s, t, r))\) ，其中
\[
\left\{\begin{array}{l}
x(s, t, r)=\frac{1}{9}(a s+4 t+8 r) \\
y(s, t, r)=\frac{1}{9}(4 s+b t-4 r) \\
z(s, t, r)=\frac{1}{9}(8 s-4 t+c r)
\end{array}\right.
\]
试讨论是否存在常数 \(a, b, c\) ，使得当 \(x=x(s, t, r), y=y(s, t, r), z=z(s, t, r)\) 时，总成立
\[
\left.\left(u_{x x}+u_{y y}+u_{z z}\right)\right|_{(x, y, z)}=\left.\left(v_{s s}+v_{t t}+v_{r r}\right)\right|_{s, t, r} .
\]
若存在，求 \(a, b, c\) 的值．

\subsection{题目6}
六．（10 分）设 \(u(x, y, z)=\frac{x^{2}}{2}+\frac{y^{2}}{3}+\frac{z^{2}}{4}\) ，求在原点处函数 \(u\) 增长最快的方向．

\subsection{题目7}
七．（10 分）证明不等式： \(\int_{0}^{\frac{\pi}{2}} \frac{\sin x}{1+x^{2}} \mathrm{~d} x<\int_{0}^{\frac{\pi}{2}} \frac{\cos x}{1+x^{2}} \mathrm{~d} x\) ．

\subsection{题目8}
八．（10 分）设函数 \(f\) 在 \([0,1]\) 上二阶可导，\(f(0)=0, f(1)=1\) ，并存在 \(c \in(0,1)\) ，使得 \(f(c)>c^{2}\) ．证明：至少存在 \(\xi \in(0,1)\) ，使得 \(f^{\prime \prime}(\xi)<2\) ．

\subsection{题目9}
九．（10 分）设 \(f\) 是 \([0,+\infty)\) 上无界的连续函数．问： \(\int_{0}^{+\infty} f(x) \mathrm{d} x\) 是否发散？给出证明或反例．

\subsection{题目10}
附加题．（10 分）设 \(\left\{f_{n}(x)\right\}\) 是 \([a, b]\) 上黎曼可积的函数列，且 \(\left\{f_{n}(x)\right\}\) 在 \([a, b]\) 上一致收敛于 \(f\) ．证明：

（1） \(f\) 在 \([a, b]\) 上可积。
（2） \(\left\{F_{n}(x)\right\}\) 在 \([a, b]\) 上一致收敛于 \(F(x)\) ，其中 \(F_{n}(x)=\int_{a}^{x} f_{n}(t) \mathrm{d} t, F(x)=\int_{0}^{x} f(t) \mathrm{d} t\) ．

\subsection{题目11}
十一．（10 分）设 \(u\) 是平面开区域 \(D\) 上的二元函数，且所有的偏导数连续．证明：\(u\) 是 \(D\) 上的调和函数，即在
\(D\) 上 \(\frac{\partial^{2} u}{\partial x^{2}}+\frac{\partial^{2} u}{\partial y^{2}} \equiv 0 \Leftrightarrow\) 对 \(D\) 内任意圆周 \(L\) ，有 \(\oint_{L} \frac{\partial u}{\partial n}=0\) ，其中 \(\frac{\partial u}{\partial n}\) 表示 \(u\) 在 \(L\) 上的外法向导数。

\subsection{题目12}
十二．（10 分）证明：存在 \(\xi \in(0, \pi)\) ，使得 \(\sum_{n=1}^{\infty} \frac{\cos n \xi}{n \ln ^{2}(n+2)}=0\) ．
\end{document}
